\scanpage{28}
\begingroup\makeatletter\tagsleft@true\makeatother
C'est un groupe unipotent, et on peut introduire
\[
Q=\underline{\operatorname{Isom}}_{\substack{\otimes,\,\text{ind. l'ident.}\\\text{sur }\operatorname{gr}^{*}}}(T_{\mathrm{DR}},T_{\mathrm{Hdg}}),\tag*{(9)}
\]
de sorte que $Q$ est un torseur à droite sous $H'$. Posons

\reversemarginpar
\margin{\figure{\begin{tikzpicture}[x=0.42cm,y=0.42cm]
\draw[->] (0.4,-3.3) .. controls (-0.55,-1.5) and (-0.05,-0.6) .. (0.55,-0.45) .. controls (1.3,-0.25) and (2.1,-0.13) .. (2.9,0.25);
\draw (1.5,-3.08) .. controls (1.12,-3.32) and (0.51,-3.05) .. (0.45,-3.3);
\draw (0.52,-3.34) .. controls (0.38,-3.92) and (0.67,-4.08) .. (1.12,-4.04);
\end{tikzpicture}}}
\normalmarginpar
\[
\begin{aligned}
\deleted{\unsure{G''|G'}}\ H''=Q\overset{H'}{\times}H=\operatorname{ad}(H')&\simeq\deleted{\operatorname{Aut}}\ \text{ss-gpe de}\\
&H''=\underline{\operatorname{Aut}}_{\unsure{\ominus|\otimes}}(T_{\mathrm{Hdg}})\ \text{respectant}\\
&\text{la filtration associée}\\
&\text{à la bigraduation,}\\
&\text{et induisant l'identité}\\
&\text{sur gradué}\\
&\text{associé --}
\end{aligned}\tag*{(1\deleted{\unsure{0|1}}\unsure{1|0})}
\]
\[
G''=Q\overset{H'}{\times}G'=\underline{\operatorname{Aut}}_{\ominus}(T_{\mathrm{Hdg}}).\tag*{(10)}
\]
Alors la bigraduation de $T_{\mathrm{Hdg}}$ définit
\[
i_1,i_2:\unsure{\mathbb{G}_{m_k}|G_{m_k}}\relbarG'\tag*{(12)}
\]
satisfaisant
\[
\varepsiloni_1(\lambda)=\varepsilon i_2(\lambda)=\lambda,\qquad \unsure{i_1|i'_1}(\lambda)i_2(\lambda)=\unsure{i|i'}(\lambda).\tag*{(13)}
\]

Ainsi, on obtient les données suivantes

\begin{enumerate}[label=\alph*)]
\item Gpe algébrique $G$ sur $\mathbb Q$, avec \diagram{\begin{tikzcd}[ampersand replacement=\&,column sep=small]\mathbb{G}_m\arrow[r,"i",no head]\& G\arrow[r,"\varepsilon",no head]\& \mathbb{G}_m\ \ldots\end{tikzcd}}
\item Groupe algébrique $G''$ sur $k$, avec \diagram{\begin{tikzcd}[ampersand replacement=\&,column sep=small]G_m\arrow[r,"i''",no head]\& G''\arrow[r,"\varepsilon''",no head]\& G_m\ \ldots\end{tikzcd}}
\item hom $i_1,i_2$ (12) satisfaisant les conditions (13)
\item Torseur \inserted{à g.} $Q$ sous le groupe \inserted{unipotent} $H''$ \inserted{$\subset G''$} défini par $(G'',\unsure{i_1|i_+})$\deleted{\unsure{le|de}}

On considère alors $\deleted{\unsure{Q\overset{H''}{\times}G''=|Q\overset{H'}{\times}G'=}}$
\[
\left\{\begin{aligned}
H'&=H''\overset{H''}{\times}Q=\operatorname{ad}_{H''}(Q)\\
\makebox[0pt][r]{\smash{\raisebox{-.8ex}{\parbox[c]{1.05em}{\figure{\tikz[x=.6em,y=.8em]{\draw[->] (.55,1.4) .. controls (-.30,1.13) and (-.27,.05) .. (.72,-.08);}}}}}\hspace{.25em}}G''&\overset{H''}{\times}Q=P\qquad\text{(tors. à g. sous $G''$)}\\
G'&=G''\overset{G''}{\times}Q:\operatorname{ad}_{G''}(Q)
\end{aligned}\right.
\]
[NB, $Q$ est trivial, donc $G'\simeq G''$ (pas canoniquement], puis
\item Torseur à g. $\unsure{P|\mathord{?}}$ sous $G'$, \deleted{avec un \unsure{isomorphisme|homomorphisme} de \unsure{$G_m\overset{G'}{\times}P$|$G'\overset{G'}{\times}P$}.}
\deleted{On considère alors $G$}
\item Un isom $G\otimes_{\mathbb Q}\unsure{K|k}\simeq G'\overset{\unsure{G'|B'}}{\times}\unsure{P^{*}|P}=\operatorname{ad}_{G'}(P)$ \margin{compatible avec les structures $i,\varepsilon$}
\item[] \deleted{\unsure{g) un|g) Un}} g) Un élément \deleted{\unsure{de|en}} $\varphi\in P(\mathbb C]$, \deleted{\illegible{1}}
\end{enumerate}
\endgroup
% archive sheet number: 27
